\subsubsection{Cylinder: Three-Dimensional Vortex Activity}\label{sec:cylinder}
Cylinder extends the volumetric evaluation to the direct numerical square-cylinder flow of Camarri et al.~\cite{Camarri:2005:Cylinder}, uniformly resampled by Weinkauf~\cite{Weinkauf:SquareCylinder,Funck:2008:Smoke}. All 102 public velocity volumes are used. We compute the three-dimensional $Q=\tfrac12(\|\Omega\|_F^2-\|S\|_F^2)$ on the $192\times64\times48$ velocity grid before retaining the author's $135\times64\times48$ downstream subdomain. The input is volumetric at every time, without spatial downsampling, smoothing, or feature-count selection. The 102 available times differ from the 508 reported in TMTM; this is a corresponding public-data experiment rather than a replication of its final derived input.

Figure~\ref{fig:cylinder} relates the P/R20 complete-sequence maps to three-dimensional source isosurfaces at frames 0, 51, and 101. A fixed $x$ reference permits comparison of downstream feature positions. The source views show the spatial distribution of positive-$Q$ structures, while the maps expose how their scalar patterns are arranged over the sequence. Persistence simplification is fixed at 2\% of each frame's range, yielding 4--75 retained extrema; it is distinct from the RG relaxation eligibility parameter. The unified evaluation tables include Cylinder with all other inputs; Supplements S9--S10 retain all five configurations and the full visual comparison.

P decreases mean NROE from 26.62\% to 12.26\% with zero measured NMLE. R2 reaches 0.249\% NROE at 0.692\% mean and 4.886\% maximum NMLE; R20 reaches 0.084\% NROE at 2.810\% mean and 14.281\% maximum NMLE. R100 has the same mean NROE as R20 but greater mean scalar cost (3.190\%). Thus, wider relaxation eligibility does not ensure a preferable measured trade-off.

For 3196 clear three-dimensional approach/separation events, errors decrease from 2141 for TMTM to 1191 for P and 605 for R20 (66.99\%, 37.27\%, and 18.93\%). Exact $x$ projection still makes 603 errors (18.87\%), identifying spatial information lost in projection. The corresponding projected-event rates use a different denominator, 2934 events: R20 makes 108 errors (3.68\%). These overlap-based geometric proxies do not establish physical vortex identities. Supplement S9 supplies the complete counts, boundary treatment, and solver records.


Figure~\ref{fig:cylinder}(d)--(f) follows a conflict at frame 86 selected by the median obstruction/gain criterion across 85 eligible frames. The source image takes the maximum $Q$ over $y$ and projects the same original three-dimensional extrema onto $(x,z)$; it does not replace the volumetric input with a planar field. P orders the marked features $A,B,C$, whereas the reference and R2 place $C$ between $A,B$. R2 lowers the frame's maximum displacement from 12.40\% to 1.17\% of $L$, at 1.003\% NMLE. The $A/B$ merge changes from 0.034060 to $-0.224935$ in $Q$ units. The local profiles use R2, while the overview maps and common primary tables use R20.
